\section{Tipo de investigación}
\label{sec:tipo_investigacion}

El estudio combinó un enfoque mixto: por un lado recogió opiniones sobre qué tan fácil y clara resultó la plataforma para los postulantes, y por otro midió datos concretos como el tiempo de espera en ventanilla, la rapidez de respuesta del servidor y si los cupos se mantuvieron sin duplicarse.

Es una investigación aplicada y tecnológica: no buscó solo plantear teoría, sino resolver un problema real del MÓDULO TECNOLÓGICO PRODUCTIVO SAN JULIÁN aplicando métodos de ingeniería de software web.

El alcance es descriptivo, propositivo y evaluativo: describe cómo funcionaba la inscripción, propone el diseño de la solución completa y luego verifica si esa solución funciona y se puede sostener, mirando tiempo, trazabilidad, consistencia de cupos y costo-beneficio.

El diseño es pre-experimental, con medición antes (2022 a 2025: registro manual y V1.0 sin recuperación de contraseña ni matrícula documental) y después (2026 con V2.0), transversal y propositivo \parencite{hernandez2014}. Sigue las cuatro fases del método científico previstas en los objetivos:

\begin{enumerate}
    \item \textbf{Fase 1 (Diagnóstico empírico):} Relevamiento de campo y análisis retrospectivo sobre las deficiencias de la plataforma preliminar de 2025 (crisis de credenciales, ruptura en la matriculación documental y asignación manual de paralelos).
    \item \textbf{Fase 2 (Diseño y modelado):} Estructuración visual de diagramas UML, modelado de la base de datos relacional y diseño de interfaces centradas en el usuario.
    \item \textbf{Fase 3 (Desarrollo técnico):} Construcción modular de la plataforma web en su Versión 2.0 bajo patrón MVC empleando Vista en HTML5/CSS3/JavaScript, Controlador en PHP 8.2/Apache y Modelo PDO--PostgreSQL 16.4.
    \item \textbf{Fase 4 (Evaluación y validación técnica, socioeducativa y de sostenibilidad):} Ejecución de pruebas funcionales de caja negra, pruebas de concurrencia y estrés, contrastación de satisfacción (OE4) y evaluación de costo-beneficio y sostenibilidad económica (OE5: inversión de 13556 Bs frente a ahorro superior a 20000 Bs por gestión) con la comunidad educativa \parencite{hernandez2014}.
\end{enumerate}

\subsection{Operacionalización de variables}

La Tabla~\ref{tab:operacionalizacion} traduce la hipótesis tecnológica en variables observables y medibles para la validación empírica en el Capítulo~7.

\begin{table}[htbp]
    \centering
    \caption{Matriz de operacionalización de la Variable Independiente y Variables Dependientes.}
    \label{tab:operacionalizacion}
    \begin{tabularx}{\textwidth}{L L L L}
        \toprule
        \textbf{Tipo} & \textbf{Variable} & \textbf{Definición operacional / Indicador} & \textbf{Unidad e instrumento} \\
        \midrule
        Independiente & V2.0 Sistema Web MTP San Julián MVC (Vista HTML5/CSS3/JS, Controlador PHP 8.2/Apache, Modelo PDO--PostgreSQL 16.4 ACID) & Plataforma integral operativa (preinscripción, recovery C.I.+fecha nacimiento, matriculación auditada, asignación por colegio) & Binaria (0/1); checklist despliegue — Cap.~4 \\
        Dependiente 1 & Tiempo de gestión del trámite & Tiempo promedio preinscripción y recovery & Minutos (3 a 5 min web vs 1 a 3 días fila; recovery $<$2 min); logs Apache + benchmarking \\
        Dependiente 2 & Tasa de recuperación autónoma & Recuperaciones exitosas sin ventanilla / intentos & Porcentaje (\%); cuestionario Likert autonomía + caja negra C.I./fecha \\
        Dependiente 3 & Trazabilidad documental & Expedientes completos auditables en BD & Porcentaje (\%) 100\%; auditoría integridad PostgreSQL \\
        Dependiente 4 & Consistencia de cupos & Ausencia de sobrecupos bajo concurrencia & Porcentaje (\%) 100\% + peticiones/min; Apache Benchmark $>$100 req/min, latencia $<$1.2 s, \texttt{FOR UPDATE} \\
        Dependiente 5 & Satisfacción usabilidad móvil & Percepción interfaz intuitiva y clara & Porcentaje (\%) $\geq$90\% (92\% esperado); Likert 1-5 a 60 estudiantes \\
        Dependiente 6 & Tasa de errores de registro & Datos inconsistentes en BD & Porcentaje (\%) $<$5\% (35\%→$<$5\%); validaciones JS + auditoría BD \\
        Dependiente 7 & Ahorro económico familiar & Diferencial pasajes/refrigerios/papelería & Bolivianos (Bs) $>$20000 por gestión; análisis costo-beneficio (inversión 13556) \\
        \bottomrule
    \end{tabularx}
    \notatabla{Elaboración propia con base en la hipótesis tecnológica del Capítulo~2 y la ficha del proyecto (Sec.~5.1). VD contrastadas en el Capítulo~7 y cerradas en el Capítulo~9.}
\end{table}
